%%% The main file. It contains definitions of basic parameters and includes all other parts.

% Meta-data of your thesis (please edit)
\input metadata.tex

% Generate metadata in XMP format for use by the pdfx package
\input xmp.tex

%% Settings for single-side (simplex) printing
% Margins: left 40mm, right 25mm, top and bottom 25mm
% (but beware, LaTeX adds 1in implicitly)
\documentclass[12pt,a4paper]{report}
\setlength\textwidth{145mm}
\setlength\textheight{247mm}
\setlength\oddsidemargin{15mm}
\setlength\evensidemargin{15mm}
\setlength\topmargin{0mm}
\setlength\headsep{0mm}
\setlength\headheight{0mm}
% \openright makes the following text appear on a right-hand page
\let\openright=\clearpage

%% Settings for two-sided (duplex) printing
% \documentclass[12pt,a4paper,twoside,openright]{report}
% \setlength\textwidth{145mm}
% \setlength\textheight{247mm}
% \setlength\oddsidemargin{14.2mm}
% \setlength\evensidemargin{0mm}
% \setlength\topmargin{0mm}
% \setlength\headsep{0mm}
% \setlength\headheight{0mm}
% \let\openright=\cleardoublepage

%% If the thesis has no printed version, symmetric margins look better
% \documentclass[12pt,a4paper]{report}
% \setlength\textwidth{145mm}
% \setlength\textheight{247mm}
% \setlength\oddsidemargin{10mm}
% \setlength\evensidemargin{10mm}
% \setlength\topmargin{0mm}
% \setlength\headsep{0mm}
% \setlength\headheight{0mm}
% \let\openright=\clearpage

%% Generate PDF/A-2u
\usepackage[a-2u]{pdfx}

%% Prefer Latin Modern fonts
\usepackage{lmodern}
% If we are not using LuaTeX, we need to set up character encoding:
\usepackage{iftex}
\ifpdftex
\usepackage[utf8]{inputenc}
\usepackage[T1]{fontenc}
\usepackage{textcomp}
\fi

%% Further useful packages (included in most LaTeX distributions)
\usepackage{amsmath}        % extensions for typesetting of math
\usepackage{amsfonts}       % math fonts
\usepackage{amsthm}         % theorems, definitions, etc.
\usepackage{bm}             % boldface symbols (\bm)
\usepackage{booktabs}       % improved horizontal lines in tables
\usepackage{caption}        % custom captions of floating objects
\usepackage{dcolumn}        % improved alignment of table columns
\usepackage{floatrow}       % custom float environments
\usepackage{graphicx}       % embedding of pictures
\usepackage{indentfirst}    % indent the first paragraph of a chapter
\usepackage[nopatch=item]{microtype}   % micro-typographic refinement
\usepackage{paralist}       % improved enumerate and itemize
\usepackage[nottoc]{tocbibind} % makes sure that bibliography and the lists
			    % of figures/tables are included in the table
			    % of contents
\usepackage{xcolor}         % typesetting in color

% The hyperref package for clickable links in PDF and also for storing
% metadata to PDF (including the table of contents).
% Most settings are pre-set by the pdfx package.
\hypersetup{unicode}
\hypersetup{breaklinks=true}

% Packages for computer science theses
\usepackage{algpseudocode}  % part of algorithmicx package
\usepackage{algorithm}
\usepackage{fancyvrb}       % improved verbatim environment
\usepackage{listings}       % pretty-printer of source code

% You might want to use cleveref for references
% \usepackage{cleveref}

% Set up formatting of bibliography (references to literature)
% Details can be adjusted in macros.tex.
%
% BEWARE: Different fields of research and different university departments
% have their own customs regarding bibliography. Consult the bibliography
% format with your supervisor.
%
% The basic format according to the ISO 690 standard with numbered references
\usepackage[natbib,style=iso-numeric,sorting=none]{biblatex}
% ISO 690 with alphanumeric references (abbreviations of authors' names)
%\usepackage[natbib,style=iso-alphabetic]{biblatex}
% ISO 690 with references Author (year)
%\usepackage[natbib,style=iso-authoryear]{biblatex}
%
% Some fields of research prefer a simple format with numbered references
% (sorting=none tells that bibliography should be listed in citation order)
%\usepackage[natbib,style=numeric,sorting=none]{biblatex}
% Numbered references, but [1,2,3,4,5] is compressed to [1-5]
%\usepackage[natbib,style=numeric-comp,sorting=none]{biblatex}
% A simple format with alphanumeric references:
%\usepackage[natbib,style=alphabetic]{biblatex}

% Load the file with bibliography entries
\addbibresource{bibliography.bib}

% Our definitions of macros (see description inside)
\input macros.tex

%%% Title page and various mandatory informational pages
\begin{document}
%%% Title page of the thesis and other mandatory pages

%%% Inscriptions at the opening page of the hard cover

% We usually do not typeset the hard cover, but if you want to do it, change \iffalse to \iftrue
\iffalse

\pagestyle{empty}
\hypersetup{pageanchor=false}
\begin{center}

\large
Charles University

\medskip

Faculty of Mathematics and Physics

\vfill

{\huge\bf\ThesisTypeTitle}

\vfill

{\huge\bf\ThesisTitle\par}

\vfill
\vfill

\hbox to \hsize{\YearSubmitted\hfil \ThesisAuthor}

\end{center}

\newpage\openright
\setcounter{page}{1}

\fi

%%% Title page of the thesis

\pagestyle{empty}
\hypersetup{pageanchor=false}
\begin{center}

\centerline{\mbox{\includegraphics[width=166mm]{img/logo-en.pdf}}}

\vspace{-8mm}
\vfill

{\bf\Large\ThesisTypeTitle}

\vfill

{\LARGE\ThesisAuthor}

\vspace{15mm}

{\LARGE\bfseries\ThesisTitle\par}

\vfill

\Department

\vfill

{
\centerline{\vbox{\halign{\hbox to 0.45\hsize{\hfil #}&\hskip 0.5em\parbox[t]{0.45\hsize}{\raggedright #}\cr
Supervisor of the \ThesisTypeName{} thesis:&\Supervisor \cr
\ifx\ThesisType\TypeRig\else
\noalign{\vspace{2mm}}
Study programme:&\StudyProgramme \cr
\fi
}}}}

\vfill

Prague \YearSubmitted

\end{center}

\newpage

%%% A page with a solemn declaration to the thesis

\openright
\hypersetup{pageanchor=true}
\vglue 0pt plus 1fill

\noindent
I declare that I carried out this \ThesisTypeName{} thesis on my own, and only with the cited
sources, literature and other professional sources.
I understand that my work relates to the rights and obligations under the Act No.~121/2000 Sb.,
the Copyright Act, as amended, in particular the fact that the Charles
University has the right to conclude a license agreement on the use of this
work as a school work pursuant to Section 60 subsection 1 of the Copyright~Act.

\vspace{10mm}

\hbox{\hbox to 0.5\hsize{%
In \hbox to 6em{\dotfill} date \hbox to 6em{\dotfill}
\hss}\hbox to 0.5\hsize{\dotfill\quad}}
\smallskip
\hbox{\hbox to 0.5\hsize{}\hbox to 0.5\hsize{\hfil Author's signature\hfil}}

\vspace{20mm}
\newpage

%%% Dedication

\openright

\noindent
\Dedication

\newpage

%%% Mandatory information page of the thesis

\openright
{\InfoPageFont

\vtop to 0.5\vsize{
\setlength\parindent{0mm}
\setlength\parskip{5mm}

Title:
\ThesisTitle

Author:
\ThesisAuthor

\DeptType:
\Department

Supervisor:
\Supervisor, \SupervisorsDepartment

Abstract:
\Abstract

Keywords:
{\def\sep{\unskip, }\ThesisKeywords}

\vfil
}

% In Czech study programmes, it is mandatory to include Czech meta-data:

\ifx\StudyLanguage\LangCS

\vtop to 0.49\vsize{
\setlength\parindent{0mm}
\setlength\parskip{5mm}

Název práce:
\ThesisTitleCS

Autor:
\ThesisAuthor

\DeptTypeCS:
\DepartmentCS

Vedoucí bakalářské práce:
\Supervisor, \SupervisorsDepartment

Abstrakt:
\AbstractCS

Klíčová slova:
{\def\sep{\unskip, }\ThesisKeywordsCS}

\vfil
}

\fi

}

\newpage

%%% Further pages will be numbered, which LaTeX automatically enables at the next chapter start


%%% A page with automatically generated table of contents of the thesis

\tableofcontents

%%% Each chapter is kept in a separate file
\chapter*{Introduction}
\addcontentsline{toc}{chapter}{Introduction}
\xxx{
	Since the creation of the World Wide Web, developers have created web applications mainly by writing and modifying text-based code.
	This approach is mainly used even today, although made easier by many different tools, languages, frameworks, and programming systems.
	//TODO
}

\section*{Problem description}
Low-code programming sytems are systems that aim to significantly reduce the amount of code needed to create certain applications \cite{Pinho_Aguiar_Amaral_2023}.
They can also help less experienced developers, by providing them with an easy-to-use interface, through which the developers can create program artifacts.

This thesis presents a novel data-driven approach for creating the client-side part of dynamic single-page applications.
We aim to create a low-code programming system based on a data-driven approach.
The data-driven approach involves first providing data to the application, which serves as the model of the application.
Then, the programming system aids the developer in creating the user interface by providing suggestions based on the structure and type of the data.
After creating the user interface, the developer can download the text-based representation of the application and modifying the behavior of these components and how they change the data itself.
This system will allow the developers to create the \emph{User Interface} of the application without the need to write a large amount of boilerplate code.

\section*{Goals of the thesis}
//TODO

\section* {Contributions and Outline}

In Chapter \ref{chap:walktrough} we showcase a creation of a user interface for a simple to-do list application.\\
In Chapter \ref{chap:design} we describe the design principles which influenced the design of the proposed programming system.\\
In Chapter \ref{chap:core-logic} we describe the type system, describe the data analysis that the system makes, and describe how the system interprets the data provided to it.\\
In Chapter \ref{chap:implementation} we describe the implementation specifics of the prototype application.\\
In Chapter \ref{chap:discussion} //TODO \\

\chapter{Background}
\label{chap:background}

\section{Low-code development systems}


\xxx{Low-code development is defined by \citet{Pinho_Aguiar_Amaral_2023} as \emph{"an approach for software development that uses tools that minimize (or eliminate) the number of lines of code written."}
	This definition accommodates a broad spectrum of tools and development practices, such as visual programming tools,
	integrated development environments with code generation capabilities, development tools powered by artificial intelligence, and generators that generate source code based on template files.
	In this thesis, we will mainly focus on visual programming tools. }

\xxx{Visual programming tools present a graphical user interface to the user. This graphical user interface serves as an abstraction over
	the development domain. It provides tools to create and modify software elements. }

\subsection{Hypercard}
\xxx{
	\emph{Hypercard} is a low-code development system created by Bill Atkinson for the Macintosh operating system. Apple released the program in 1987 at the
	Macworld exposition in Boston~\cite{hyper_release}. Apple developed and maintained the program until 1998. The program became very popular, and similar programs and clones of Hypercard were created after the discontinuation.
}

\xxx{The following summary of functionality is based on a book written by~\citet{goodman_hypertext}.
	The fundamental elements that the user creates are called \emph{cards}. Cards can hold data as text, have custom graphics, contain buttons, and implement custom behavior.
	Users can implement the custom behavior using a scripting language called \emph{HyperTalk}. Then, users can group cards into \emph{stacks}. A stack is a collection of cards with the same type of information.
	The program saves a stack as a single file to the disk. Finally, users can distribute and modify these stacks.
}

\begin{figure}
	\centering
	\includegraphics[width=.7\linewidth]{img/hypercard−menu.pdf}
	\caption{Selection of a user level in the Hypercard preferences menu.}
	\label{fig:f}
\end{figure}

\xxx{
	The program offers users a choice of a user level, as seen in figure \ref{fig:f}. The program changes the user interface's capabilities based on the selected user level.
	The higher the selected user level, the more options the program enables and displays. A selected level enables all previous-level options alongside other options and functionality.
	Users can choose from five different user levels:
	\begin{enumerate}
		\item Browsing - enables no modification of cards and is mainly intended to be used to view the different cards inside a stack.
		\item Typing - enables inserting and modifying text data inside the cards.
		\item Painting - enables the creation of custom graphics inside the cards. The program provides different graphical options and tools.
		\item Authoring - adds the ability to add buttons and fields. This way, the users can add card functionality without writing code.
		\item Scripting - adds the ability to use the HyperTalk scripting language to modify the behavior of the different elements inside the cards.
	\end{enumerate}
}

\xxx{
	The program became popular with users at many different skill levels thanks to separating more advanced features from the basic ones.
	Less advanced users can learn to use the interface more easily and gradually increase their proficiency with the program without being overwhelmed by options.
	On the other hand, advanced users can set the highest user level and use the program fully from the start. }


\subsection{Darklang}
\xxx{\emph{Darklang}\footnote{\url{https://darklang.com/} [visited on 2024-04-22][online]} is a cloud-based low-code programming system for building web application backends created by Paul Biggar.
	The original system's development was discontinued in 2023 and renamed Darklang-classic.
}

\xxx{The following summary follows the official Darklang~documentation\footnote{\url{https://docs.darklang.com/} [visited on 2024-04-22][online]}.
	The Darklang project consists of a low-code editor and a programming language. It also provides tools for creating persistent storage and for deployment.
	The editor provides the main functionality of the service. It consists of a canvas that displays the created application elements as draggable boxes.
	Each box provides options to modify the element. These options include different menus and input fields.
	Some elements can be modified using the Darklang programming language to implement custom behavior.
	The user can create program elements of the following categories:
	\begin{itemize}
		\item HTTP handlers - definition of API endpoints
		\item Persistent storage - database creation and modification
		\item Workers - processing of background tasks
		\item Cron jobs - scheduled jobs with custom behavior
		\item REPL - a general-purpose programmable element
		\item Functions - a reusable element with custom behavior
  \end{itemize}
	Through the combination of these elements, the users can create application backends of varying complexity.
	After the user creates and modifies a particular element, the system automatically redeploys the entire application.
	This allows for an interactive style of development where changes made to the system are almost immediately visible.}

\xxx{Users can implement custom functionality by using the \emph{Darklang language}.
	It is a statically-typed \emph{ML-style} language similar to \emph{OCaml}\footnote{\url{https://ocaml.org/} [visited on 2024-04-23][online]}.
	Some of the features of the language include:
	\begin{itemize}
		\item The pipe operator.
		\item Functions for working with collections, such as the \emph{map} and \emph{fold} functions.
		\item Pattern matching.
		\item Implicit returns.
	\end{itemize}
}

\subsection{Output-directed programming}
\xxx{
	\emph{Output-directed programming} is described by \citet{sketch-and-sketch} as
	a programming paradigm where users construct plain text programs using mouse-based operations on the program's output.
	The text representation of the program serves as the main representation.
	This means that the user can also directly change the text of the program, which will result in a change in output.
	This allows the user to make changes that are difficult or impossible to make using the direct manipulation of the program's output.
	This bimodal approach can be seen as an extension of the standard programming paradigm where the manipulation of the text-based program
	representation is the only way to change the program's output.
}

\xxx{One example of a programming system that demonstrates this paradigm is called \emph{Sketch-and-Sketch} firstly presented by \citet{output-directed-programming}.
	The system allows the creation and modification of Scalable Vector Graphics using direct manipulation of the program's SVG output.
}

\section{Web development tools}
\subsection{Elm}
\subsection{Web development using F\# }

\chapter{Implementation analysis}

\section{}

\chapter{Design principles}
\label{chap:design}

\chapter{Core logic}
\label{chap:corelogic}

This chapter describes the core logic of the programming system.
Firstly, we describe the hole-based approach applied to the creation of elements.
After that, we describe the type system, mainly the internal representation of the UI elements and necessary operations on these types.
Then, we describe mapping the input data types to specific types representing the UI elements.
We present two ways of combining the created UI elements with the corresponding data.
Lastly, we describe one possible way of displaying a preview of the created UI elements and the modification menus.

\section{Hole-based approach}
\label{sec:hole-based}
The primary motivation of this approach is to allow \emph{incremental creation} of UI elements.
The elements are created based on the provided data and its structure.
Each element corresponds to a data value.
We need the ability to represent data values that have yet to be used to create UI elements.
To address this need, we define a type named \emph{Hole} as follows:
\begin{defn}[Hole type definition]
	A Hole is a placeholder type representing a UI element that is yet to be created despite its existing corresponding data. It serves as a temporary stand-in during the incremental construction of a user interface.
\end{defn}

We then define the incremental creation process as a \emph{sequence of discrete opreations}.
Each operation is either a \emph{modification} of an existing UI element or a \emph{replacement} of a Hole element with a new UI element based on the corresponding data.
This discrete approach allows the system to perform different tasks after each operation.
These tasks include updating the UI element preview, analyzing the input data, or providing new modification menus and options.
The system could also analyze the existing UI elements and other operations.


\section{Type system and Operations}

\subsection{JSON type}
\label{sub:json}
The creation process starts by providing \emph{JSON} data to the system.
In order to use this data, we parse it and create an internal representation.
We define a discriminated union type named \emph{JSON}.
This type is the foundation for the Abstract Syntax Tree (AST) representation of the JSON data.
The type is defined as a as follows:
\begin{listing}[H]
	\caption {JSON type}
	\begin{lstlisting}
  type Json =
    | JNumber of float
    | JString of string
    | JBool of bool
    | JNull
    | JArray of Json list
    | JObject of Map<string, Json>
  \end{lstlisting}
\end{listing}

Each node represents a JSON value in the input data.
The types of nodes can be divided into two categories:
\begin{itemize}
	\item {\textbf{Collections:} The first category contains types representing a \emph{collection} of other values. This category includes the types \emph{JObject} and \emph{JArray}.
	      We define the two collection types as follows:
	      \begin{itemize}
		      \item \textbf{JObject:} It is based on the JSON Object type and represents a collection of different JSON types. The original ordering of the inner elements is \emph{ignored}.
		      \item \textbf{JArray:} It is based on the JSON List type and represents a collection of JSON values of the \emph{same type and structure}. The original ordering of the inner elements is \emph{preserved}.
	      \end{itemize}
	      }
	\item \textbf{Primitives:} The second category contains types representing data primitives such as numerical values, boolean values,
	      a string of text, or the null value.
\end{itemize}

\subsection{RenderingCode type}
To facilitate the creation of UI elements based on the JSON type, we define a type named \emph{RenderingCode}.
The \emph{RenderingCode} is a discriminated union type used to represent the UI elements.
Similarly to the JSON type, each case represents a type of an HTML element with a corresponding mapping to the JSON type.
This type is the foundation for the \emph{Abstract Syntax Tree} representation of UI elements.
\begin{listing}[H]
	\caption {RenderingCode type}
	\begin{lstlisting}
  type RenderingCode =
    | HtmlElement of
        tag: Tag *
        attrs: Attributes *
        innerValue: InnerValue
    | HtmlList of
        listType: ListType *
        elements: RenderingCode list
    | HtmlObject of
      objectType : ObjectType *
      elements: RenderingCode list
    | Hole
  \end{lstlisting}
\end{listing}

The case types are defined as follows:
\begin{itemize}
	\item \textbf{HtmlElement:} Represents a single HTML element. It consists of:
	      \begin{itemize}
		      \item \emph{Tag}: Denotes the HTML tag (e.g., div, p, span).

		      \item \emph{InnerValue}: Represents the content of the element, defined as:
		            \begin{lstlisting}
type InnerValue =
  | Data 
  | Empty
  | Constant of string  
        \end{lstlisting}
		            Where \emph{Data} indicates dynamic content from JSON, \emph{Empty} represents an empty element, and \emph{Constant} holds static, developer-defined content.
	      \end{itemize}	\item \textbf{HtmlList:} Represents a collection of similar HTML elements corresponding to JSON arrays. It includes:
	      \begin{itemize}
		      \item \emph{ListType}: Specifies the list type (e.g., unordered, ordered, table).
		      \item \emph{itemCodes}: A list of RenderingCode elements representing list items.
		      \item {\emph{Attributes}: A collection of key-value pairs representing HTML attributes.
		            The Attribute type is defined as follows:
		            \begin{lstlisting}
type Attribute = string * InnerValue
              \end{lstlisting}
		            It consists of a key of a type string and an value of a type InnerValue. This allows selecting the value of an Attribute as a value from the input data.
		            }

	      \end{itemize}
	\item \textbf{HtmlObject:} Represents a collection of diverse HTML elements derived from JSON object. It consists of:
	      \begin{itemize}
		      \item \emph{ObjectType}: Defines the type of the object (e.g., div, span, section).
		      \item \emph{Elements}: A list of RenderingCode elements representing the object's contents.
	      \end{itemize}
	\item {\textbf{Hole:} Represents a placeholder in the UI element structure as defined in Section~\ref{sec:hole-based}, allowing for dynamic element creation.
	      The Hole is a discriminated union type
	      }


	\item  \xxx{//TODO: CustomWrapper and Javascript}
\end{itemize}


\subsection{Data Mapping: JSON to RenderingCode}
The process of transforming JSON data into UI elements requires systematic mapping between JSON and RenderingCode types.
This mapping forms the core of our UI generation system, allowing us to convert arbitrary JSON input into a structured representation of UI elements.
Each JSON type has a corresponding representation in RenderingCode:
\begin{itemize}
	\item JObject maps to HtmlObject
	\item JArray maps to HtmlList
	\item JNull, JString, JNumber, and JBool map to HtmlElement
\end{itemize}
Th

We define a mapping function named \emph{recognizeJson} as follows:
\begin{listing}[H]
	\caption {RenderingCode type}
	\begin{lstlisting}
let recognizeJson (json: Json) =
    match json with
    | JArray array -> 
        //create a hole for each element of the array
        HtmlList containing the holes as elements
    | JObject obj ->
        //create hole for each element of the object 
        HtmlObject containing the holes as elements
    | JNull -> Hole(UnNamed)
    | _ -> HtmlElement 
  \end{lstlisting}
\end{listing}


\section{Creation process}



\section{Referencing input data}
\subsection{Structual referencing}
\subsection{Selector based approach}


\chapter{Results}
\chapter*{Conclusion}
\addcontentsline{toc}{chapter}{Conclusion}


%%% Bibliography
%%% Bibliography (literature used as a source)
%%%
%%% We employ biblatex to construct the bibliography. It processes
%%% citations in the text (e.g., the \cite{...} macro) and looks up
%%% relevant entries in the bibliography.bib file.
%%%
%%% See also biblatex settings in thesis.tex.

%%% Generate the bibliography. Beware that if you cited no works,
%%% the empty list will be omitted completely.

% We let bibliography items stick out of the right margin a little
\def\bibfont{\hfuzz=2pt}

\printbibliography[heading=bibintoc]

%%% If case you prefer to write the bibliography manually (without biblatex),
%%% you can use the following. Please follow the ISO 690 standard and
%%% citation conventions of your field of research.

% \begin{thebibliography}{99}
%
% \bibitem{lamport94}
%   {\sc Lamport,} Leslie.
%   \emph{\LaTeX: A Document Preparation System}.
%   2nd edition.
%   Massachusetts: Addison Wesley, 1994.
%   ISBN 0-201-52983-1.
%
% \end{thebibliography}


%%% Figures used in the thesis (consider if this is needed)
\listoffigures

%%% Tables used in the thesis (consider if this is needed)
%%% In mathematical theses, it could be better to move the list of tables to the beginning of the thesis.
\listoftables

%%% Abbreviations used in the thesis, if any, including their explanation
%%% In mathematical theses, it could be better to move the list of abbreviations to the beginning of the thesis.
\chapwithtoc{List of Abbreviations}

%%% Doctoral theses must contain a list of author's publications
\ifx\ThesisType\TypePhD
\chapwithtoc{List of Publications}
\fi

%%% Attachments to the thesis, if any. Each attachment must be referred to
%%% at least once from the text of the thesis. Attachments are numbered.
%%%
%%% The printed version should preferably contain attachments, which can be
%%% read (additional tables and charts, supplementary text, examples of
%%% program output, etc.). The electronic version is more suited for attachments
%%% which will likely be used in an electronic form rather than read (program
%%% source code, data files, interactive charts, etc.). Electronic attachments
%%% should be uploaded to SIS. Allowed file formats are specified in provision
%%% of the rector no. 72/2017. Exceptions can be approved by faculty's coordinator.
\appendix
\chapter{Attachments}

\section{First Attachment}

\end{document}
